
\documentclass{article}
\usepackage{colortbl}
\usepackage{makecell}
\usepackage{multirow}
\usepackage{supertabular}

\begin{document}

\newcounter{utterance}

\centering \large Interaction Transcript for game `hot\_air\_balloon', experiment `air\_balloon\_survival\_en\_complexity\_easy', episode 0 with deepseek{-}chat{-}v3.1{-}t1.0.
\vspace{24pt}

{ \footnotesize  \setcounter{utterance}{1}
\setlength{\tabcolsep}{0pt}
\begin{supertabular}{c@{$\;$}|p{.15\linewidth}@{}p{.15\linewidth}p{.15\linewidth}p{.15\linewidth}p{.15\linewidth}p{.15\linewidth}}
    \# & $\;$A & \multicolumn{4}{c}{Game Master} & $\;\:$B\\
    \hline

    \theutterance \stepcounter{utterance}  
    & & \multicolumn{4}{p{0.6\linewidth}}{
        \cellcolor[rgb]{0.9,0.9,0.9}{
            \makecell[{{p{\linewidth}}}]{
                \texttt{\tiny{[P1$\langle$GM]}}
                \texttt{You are participating in a collaborative negotiation game.} \\
\\ 
\texttt{Together with another participant, you must agree on a single set of items that will be kept. Each of you has your own view of how much each item matters to you (importance). You do not know how the other participant values the items. Additionally, you are given the effort each item demands.} \\
\texttt{You may only agree on a set if the total effort of the selected items does not exceed a shared limit:} \\
\\ 
\texttt{LIMIT = 7605} \\
\\ 
\texttt{Here are the individual item effort values:} \\
\\ 
\texttt{Item effort = \{"B21": 369, "B00": 951, "A08": 156, "A27": 266, "B94": 581, "B53": 287, "A23": 866, "C65": 180, "A89": 800, "B67": 739, "B45": 637, "C34": 86, "C12": 749, "A05": 370, "B60": 345, "C92": 951, "B69": 147, "A84": 265, "C06": 262, "B12": 259, "A16": 358, "B85": 394, "C35": 286, "B10": 580, "C15": 479, "B09": 14, "B16": 153, "C81": 977, "A58": 134, "C24": 982, "A30": 259, "A66": 232, "B75": 202, "B46": 73, "C70": 821\}} \\
\\ 
\texttt{Here is your personal view on the importance of each item:} \\
\\ 
\texttt{Item importance values = \{"B21": 138, "B00": 583, "A08": 868, "A27": 822, "B94": 783, "B53": 65, "A23": 262, "C65": 121, "A89": 508, "B67": 780, "B45": 461, "C34": 484, "C12": 668, "A05": 389, "B60": 808, "C92": 215, "B69": 97, "A84": 500, "C06": 30, "B12": 915, "A16": 856, "B85": 400, "C35": 444, "B10": 623, "C15": 781, "B09": 786, "B16": 3, "C81": 713, "A58": 457, "C24": 273, "A30": 739, "A66": 822, "B75": 235, "B46": 606, "C70": 968\}} \\
\\ 
\texttt{Goal:} \\
\\ 
\texttt{Your goal is to negotiate a shared set of items that benefits you as much as possible (i.e., maximizes total importance to YOU), while staying within the LIMIT. You are not required to make a PROPOSAL in every message {-} you can simply negotiate as well. All tactics are allowed!} \\
\\ 
\texttt{Interaction Protocol:} \\
\\ 
\texttt{You may only use the following structured formats in a message:} \\
\\ 
\texttt{PROPOSAL: \{'A', 'B', 'C', …\}} \\
\texttt{Propose keeping exactly those items.} \\
\texttt{REFUSE: \{'A', 'B', 'C', …\}} \\
\texttt{Explicitly reject opponent's proposal.} \\
\texttt{ARGUMENT: \{'...'\}} \\
\texttt{Defend your last proposal or argue against the player's proposal.} \\
\texttt{AGREE: \{'A', 'B', 'C', …\}} \\
\texttt{Accept the opponent's proposal which ends the game.} \\
\texttt{STRATEGIC REASONING: \{'...'\}} \\
\texttt{	Describe your strategic reasoning or anticipation explaining your choice of action. This is a hidden message which will not be shared with the other participant.} \\
\\ 
\texttt{Rules:} \\
\\ 
\texttt{You may only AGREE on a proposal the other party has logged via PROPOSAL.} \\
\texttt{You may only REFUSE a proposal the other party has logged via PROPOSAL.} \\
\texttt{Total effort of any PROPOSAL or AGREE set must be ≤ LIMIT.} \\
\texttt{Do NOT reveal your hidden importance scores.} \\
\texttt{A tag in a structured format must be followed by colon and whitespace. The argument must be a python set containing 0 or more strings.} \\
\texttt{So, it must be of the form TAG: \{...\}} \\
\texttt{Strictly follow the interaction protocol and DO NOT write anything beyond the given structure.} \\
\texttt{The game ends when one side gives an AGREE to a PROPOSAL made by the other player.} \\
\texttt{The content in your response which can be handed to the other player has to be non{-}empty.} \\
\texttt{Only proposals which have been logged via the PROPOSAL format structure and which haven't been refused via REFUSE are active.} \\
\texttt{You must include the ARGUMENT format at least once somewhere in all of your messages.} \\
\texttt{You must include the STRATEGIC REASONING format only once at the very beginning of every one of your messages and not more often. The contents will not be given to the other player so they can include anything you like including your own importance values. Here you should reason step by step to come up with you next move.} \\
            }
        }
    }
    & \\ \\

    \theutterance \stepcounter{utterance}  
    & \multicolumn{4}{p{0.6\linewidth}}{
        \cellcolor[rgb]{0.9,0.9,0.9}{
            \makecell[{{p{\linewidth}}}]{
                \texttt{\tiny{[P1$\rangle$GM]}}
                \texttt{STRATEGIC REASONING: \{'I need to maximize my total importance while staying under 7605 effort. First, I\textbackslash{}'ll calculate the total effort of all items to see the scale. Total effort sum is 369+951+156+266+581+287+866+180+800+739+637+86+749+370+345+951+147+265+262+259+358+394+286+580+479+14+153+977+134+982+259+232+202+73+821 = 17221, which is way above the limit, so we need to select a subset. My top importance items are C70:968, B12:915, A08:868, A16:856, A27:822, A66:822, B60:808, B09:786, B94:783, C15:781, B67:780, A30:739, C81:713, C12:668, B10:623, B00:583, B46:606, A89:508, A84:500, C34:484, A58:457, B45:461, C35:444, B85:400, A05:389, B21:138, C65:121, B69:97, C92:215, B75:235, C24:273, A23:262, B53:65, C06:30, B16:3. But I must consider effort: high importance items with low effort are ideal. I should prioritize items with high importance/effort ratio. Calculating ratios: C34:484/86=5.63, B09:786/14=56.14, B46:606/73=8.30, A58:457/134=3.41, B69:97/147=0.66, C65:121/180=0.67, B16:3/153=0.02, C06:30/262=0.11, B53:65/287=0.23, B21:138/369=0.37, A05:389/370=1.05, B75:235/202=1.16, A30:739/259=2.85, B12:915/259=3.53, A66:822/232=3.54, A84:500/265=1.89, C35:444/286=1.55, A27:822/266=3.09, B60:808/345=2.34, B10:623/580=1.07, C15:781/479=1.63, B45:461/637=0.72, B67:780/739=1.06, A08:868/156=5.56, B94:783/581=1.35, A16:856/358=2.39, C12:668/749=0.89, B00:583/951=0.61, A89:508/800=0.64, B85:400/394=1.02, C70:968/821=1.18, C81:713/977=0.73, C92:215/951=0.23, A23:262/866=0.30, C24:273/982=0.28. So best ratios are B09, C34, A08, B46, A58, etc. But I need to consider the other player\textbackslash{}'s perspective; they might value different items. Since I don\textbackslash{}'t know their importance, I should start with a proposal that includes many high{-}ratio items for me, but also leaves room for negotiation. I should aim for a set with total importance around 8000{-}9000? But effort must be under 7605. Let me draft a proposal: include all items with ratio >1.5 or so, but avoid very high effort items with low ratio. For example, exclude C81 (977eff, 0.73), C92 (951,0.23), A23 (866,0.30), C24 (982,0.28), B00 (951,0.61), A89 (800,0.64), B45 (637,0.72), C12 (749,0.89). But include B09 (14eff,786imp), C34 (86,484), A08 (156,868), B46 (73,606), A58 (134,457), A30 (259,739), B12 (259,915), A66 (232,822), A27 (266,822), B60 (345,808), C15 (479,781), B67 (739,780), A16 (358,856), B10 (580,623), B85 (394,400), C70 (821,968), etc. Calculate effort: B09:14, C34:86, A08:156, B46:73, A58:134, A30:259, B12:259, A66:232, A27:266, B60:345, C15:479, B67:739, A16:358, B10:580, B85:394, C70:821. Sum: 14+86=100; +156=256; +73=329; +134=463; +259=722; +259=981; +232=1213; +266=1479; +345=1824; +479=2303; +739=3042; +358=3400; +580=3980; +394=4374; +821=5195. Well under 7605. I can add more. Add B75 (202,235 ratio1.16): 5195+202=5397. Add C35 (286,444 ratio1.55): 5397+286=5683. Add A84 (265,500 ratio1.89): 5683+265=5948. Add A05 (370,389 ratio1.05): 5948+370=6318. Add B21 (369,138 ratio0.37) not good. Add C65 (180,121 ratio0.67) not. Add B53 (287,65) no. Add B94 (581,783 ratio1.35): 6318+581=6899. Add C12 (749,668 ratio0.89): 6899+749=7648, which is over 7605? 7648{-}7605=43 over. So skip C12. Instead, add something smaller: B69 (147,97) no. Or C06 (262,30) no. Or B16 (153,3) no. So with B94, total 6899. Then add nothing? But I have room: 7605{-}6899=706. I can add A89? But A89 has 800eff, too big. B00? 951 too big. C81? 977 too big. C92? 951 too big. A23? 866 too big. C24? 982 too big. So maybe add multiple small ones: like B75 already added, C35 added, A84 added, A05 added. So my current set has effort 6899. I can add B85 already added. Or add C70 already added. Actually, I have only 17 items. I can add more with good ratio but not too high effort. For example, add C65? 180eff, but ratio low. Not worth. Better to keep effort low for negotiation. So my initial proposal will be: \{B09, C34, A08, B46, A58, A30, B12, A66, A27, B60, C15, B67, A16, B10, B85, C70, B75, C35, A84, A05, B94\}. This has effort 6899 and my importance: B09:786, C34:484, A08:868, B46:606, A58:457, A30:739, B12:915, A66:822, A27:822, B60:808, C15:781, B67:780, A16:856, B10:623, B85:400, C70:968, B75:235, C35:444, A84:500, A05:389, B94:783. Sum imp: let\textbackslash{}'s calculate: 786+484=1270; +868=2138; +606=2744; +457=3201; +739=3940; +915=4855; +822=5677; +822=6499; +808=7307; +781=8088; +780=8868; +856=9724; +623=10347; +400=10747; +968=11715; +235=11950; +444=12394; +500=12894; +389=13283; +783=14066. This is very high for me, but effort is only 6899, so I have room to concede items that are less important to me but might be important to the other. This is a good starting point.'\}} \\
\\ 
\texttt{PROPOSAL: \{'B09', 'C34', 'A08', 'B46', 'A58', 'A30', 'B12', 'A66', 'A27', 'B60', 'C15', 'B67', 'A16', 'B10', 'B85', 'C70', 'B75', 'C35', 'A84', 'A05', 'B94'\}} \\
\texttt{ARGUMENT: \{'This proposal includes a diverse set of items with a total effort of 6899, well under our limit of 7605. It maximizes value across many categories while leaving room for additional items you might favor. I believe this is a strong starting point for discussion.'\}} \\
            }
        }
    }
    & & \\ \\

    \theutterance \stepcounter{utterance}  
    & & \multicolumn{4}{p{0.6\linewidth}}{
        \cellcolor[rgb]{0.9,0.9,0.9}{
            \makecell[{{p{\linewidth}}}]{
                \texttt{\tiny{[GM$\rangle$P2]}}
                \texttt{You are participating in a collaborative negotiation game.} \\
\\ 
\texttt{Together with another participant, you must agree on a single set of items that will be kept. Each of you has your own view of how much each item matters to you (importance). You do not know how the other participant values the items. Additionally, you are given the effort each item demands.} \\
\texttt{You may only agree on a set if the total effort of the selected items does not exceed a shared limit:} \\
\\ 
\texttt{LIMIT = 7605} \\
\\ 
\texttt{Here are the individual item effort values:} \\
\\ 
\texttt{Item effort = \{"B21": 369, "B00": 951, "A08": 156, "A27": 266, "B94": 581, "B53": 287, "A23": 866, "C65": 180, "A89": 800, "B67": 739, "B45": 637, "C34": 86, "C12": 749, "A05": 370, "B60": 345, "C92": 951, "B69": 147, "A84": 265, "C06": 262, "B12": 259, "A16": 358, "B85": 394, "C35": 286, "B10": 580, "C15": 479, "B09": 14, "B16": 153, "C81": 977, "A58": 134, "C24": 982, "A30": 259, "A66": 232, "B75": 202, "B46": 73, "C70": 821\}} \\
\\ 
\texttt{Here is your personal view on the importance of each item:} \\
\\ 
\texttt{Item importance values = \{"B21": 138, "B00": 583, "A08": 868, "A27": 822, "B94": 783, "B53": 65, "A23": 262, "C65": 121, "A89": 508, "B67": 780, "B45": 461, "C34": 484, "C12": 668, "A05": 389, "B60": 808, "C92": 215, "B69": 97, "A84": 500, "C06": 30, "B12": 915, "A16": 856, "B85": 400, "C35": 444, "B10": 623, "C15": 781, "B09": 786, "B16": 3, "C81": 713, "A58": 457, "C24": 273, "A30": 739, "A66": 822, "B75": 235, "B46": 606, "C70": 968\}} \\
\\ 
\texttt{Goal:} \\
\\ 
\texttt{Your goal is to negotiate a shared set of items that benefits you as much as possible (i.e., maximizes total importance to YOU), while staying within the LIMIT. You are not required to make a PROPOSAL in every message {-} you can simply negotiate as well. All tactics are allowed!} \\
\\ 
\texttt{Interaction Protocol:} \\
\\ 
\texttt{You may only use the following structured formats in a message:} \\
\\ 
\texttt{PROPOSAL: \{'A', 'B', 'C', …\}} \\
\texttt{Propose keeping exactly those items.} \\
\texttt{REFUSE: \{'A', 'B', 'C', …\}} \\
\texttt{Explicitly reject opponent's proposal.} \\
\texttt{ARGUMENT: \{'...'\}} \\
\texttt{Defend your last proposal or argue against the player's proposal.} \\
\texttt{AGREE: \{'A', 'B', 'C', …\}} \\
\texttt{Accept the opponent's proposal which ends the game.} \\
\texttt{STRATEGIC REASONING: \{'...'\}} \\
\texttt{	Describe your strategic reasoning or anticipation explaining your choice of action. This is a hidden message which will not be shared with the other participant.} \\
\\ 
\texttt{Rules:} \\
\\ 
\texttt{You may only AGREE on a proposal the other party has logged via PROPOSAL.} \\
\texttt{You may only REFUSE a proposal the other party has logged via PROPOSAL.} \\
\texttt{Total effort of any PROPOSAL or AGREE set must be ≤ LIMIT.} \\
\texttt{Do NOT reveal your hidden importance scores.} \\
\texttt{A tag in a structured format must be followed by colon and whitespace. The argument must be a python set containing 0 or more strings.} \\
\texttt{So, it must be of the form TAG: \{...\}} \\
\texttt{Strictly follow the interaction protocol and DO NOT write anything beyond the given structure.} \\
\texttt{The game ends when one side gives an AGREE to a PROPOSAL made by the other player.} \\
\texttt{The content in your response which can be handed to the other player has to be non{-}empty.} \\
\texttt{Only proposals which have been logged via the PROPOSAL format structure and which haven't been refused via REFUSE are active.} \\
\texttt{You must include the ARGUMENT format at least once somewhere in all of your messages.} \\
\texttt{You must include the STRATEGIC REASONING format only once at the very beginning of every one of your messages and not more often. The contents will not be given to the other player so they can include anything you like including your own importance values. Here you should reason step by step to come up with you next move.} \\
\\ 
\texttt{You will now receive the first message of the other player.} \\
\\ 
\texttt{PROPOSAL: \{'B09', 'C34', 'A08', 'B46', 'A58', 'A30', 'B12', 'A66', 'A27', 'B60', 'C15', 'B67', 'A16', 'B10', 'B85', 'C70', 'B75', 'C35', 'A84', 'A05', 'B94'\}} \\
\texttt{ARGUMENT: \{'This proposal includes a diverse set of items with a total effort of 6899, well under our limit of 7605. It maximizes value across many categories while leaving room for additional items you might favor. I believe this is a strong starting point for discussion.'\}} \\
            }
        }
    }
    & \\ \\

    \theutterance \stepcounter{utterance}  
    & & & \multicolumn{4}{p{0.6\linewidth}}{
        \cellcolor[rgb]{0.9,0.9,0.9}{
            \makecell[{{p{\linewidth}}}]{
                \texttt{\tiny{[GM$\langle$P2]}}
                \texttt{STRATEGIC REASONING: \{'I need to evaluate the opponent's proposal from my own importance perspective. First, I calculate the total effort: 369(B21) is missing, 951(B00) missing, 156(A08) included, 266(A27) included, 581(B94) included, 287(B53) missing, 866(A23) missing, 180(C65) missing, 800(A89) missing, 739(B67) included, 637(B45) missing, 86(C34) included, 749(C12) missing, 370(A05) included, 345(B60) included, 951(C92) missing, 147(B69) missing, 265(A84) included, 262(C06) missing, 259(B12) included, 358(A16) included, 394(B85) included, 286(C35) included, 580(B10) included, 479(C15) included, 14(B09) included, 153(B16) missing, 977(C81) missing, 134(A58) included, 982(C24) missing, 259(A30) included, 232(A66) included, 202(B75) included, 73(B46) included, 821(C70) included. My calculation: (156+266+581+739+86+370+345+265+259+358+394+286+580+479+14+134+259+232+202+73+821) = 6899. This matches their stated effort. Now, I calculate my total importance for this set: A08(868), A27(822), B94(783), B67(780), C34(484), A05(389), B60(808), A84(500), B12(915), A16(856), B85(400), C35(444), B10(623), C15(781), B09(786), A58(457), A30(739), A66(822), B75(235), B46(606), C70(968). Sum: 868+822=1690, +783=2473, +780=3253, +484=3737, +389=4126, +808=4934, +500=5434, +915=6349, +856=7205, +400=7605, +444=8049, +623=8672, +781=9453, +786=10239, +457=10696, +739=11435, +822=12257, +235=12492, +606=13098, +968=14066. This gives me a total importance of 14066. I should compare this to my ideal set. I have high{-}value items like C70(968), B12(915), A08(868), A16(856), A27(822), A66(822), B60(808), B94(783), C15(781), B67(780), B09(786) which are mostly included. But I'm missing some high{-}importance items like B00(583), A89(508), C12(668), B45(461), A30(739 is included), B10(623 is included), B46(606 is included). The biggest misses for me are C12(668) and A89(508). I have low{-}importance items included like B75(235) and C35(444) which I might be willing to trade. The proposal is good for me, but I might be able to improve it by adding C12(668, effort 749) and A89(508, effort 800) which would add 1549 effort, making total 6899+1549=8448 which is over the limit of 7605. So I need to remove some items to make room. I can remove lower{-}importance items like B75(235, effort 202), C35(444, effort 286), B85(400, effort 394), which together free up 202+286+394=882 effort. This would allow me to add C12(749) and A89(800) with net effort change +749+800{-}882=667, so new total 6899+667=7566 which is under 7605. This would increase my importance by +668+508{-}235{-}444{-}400=97 (from 14066 to 14163). Alternatively, I could propose adding just C12(668,749) and remove B75(235,202) and C35(444,286), net effort +749{-}202{-}286=261, new total 6899+261=7160, importance change +668{-}235{-}444={-}11, which is worse. Or add A89(508,800) and remove B75(235,202) and C35(444,286), net effort +800{-}202{-}286=312, total 6899+312=7211, importance change +508{-}235{-}444={-}171, worse. So the best swap is to add both C12 and A89 while removing B75, C35, and B85. But I should not reveal my full strategy immediately. I should start by acknowledging the proposal and suggesting a modification to include some items I value highly, without specifying which ones exactly. I'll propose a modified set that includes my high{-}value items and see the reaction.'\}} \\
\\ 
\texttt{PROPOSAL: \{'B09', 'C34', 'A08', 'B46', 'A58', 'A30', 'B12', 'A66', 'A27', 'B60', 'C15', 'B67', 'A16', 'B10', 'B94', 'C12', 'A89', 'C70', 'A84', 'A05', 'B85'\}} \\
\texttt{ARGUMENT: \{'Thank you for the proposal. It is a good starting point, but I believe we can improve the set by including items C12 and A89, which offer high value, and adjusting by removing some lower{-}impact items to stay within the limit. This modified set maintains a strong foundation while incorporating additional high{-}value elements.'\}} \\
            }
        }
    }
     \\ \\

    \theutterance \stepcounter{utterance}  
    & & \multicolumn{4}{p{0.6\linewidth}}{
        \cellcolor[rgb]{0.9,0.9,0.9}{
            \makecell[{{p{\linewidth}}}]{
                \texttt{\tiny{[P1$\langle$GM]}}
                \texttt{PROPOSAL: \{'B09', 'C34', 'A08', 'B46', 'A58', 'A30', 'B12', 'A66', 'A27', 'B60', 'C15', 'B67', 'A16', 'B10', 'B94', 'C12', 'A89', 'C70', 'A84', 'A05', 'B85'\}} \\
\texttt{ARGUMENT: \{'Thank you for the proposal. It is a good starting point, but I believe we can improve the set by including items C12 and A89, which offer high value, and adjusting by removing some lower{-}impact items to stay within the limit. This modified set maintains a strong foundation while incorporating additional high{-}value elements.'\}} \\
            }
        }
    }
    & \\ \\

    \theutterance \stepcounter{utterance}  
    & \multicolumn{4}{p{0.6\linewidth}}{
        \cellcolor[rgb]{0.9,0.9,0.9}{
            \makecell[{{p{\linewidth}}}]{
                \texttt{\tiny{[P1$\rangle$GM]}}
                \texttt{STRATEGIC REASONING: \{'The opponent has proposed a set very similar to mine, but they added C12 (749 effort) and A89 (800 effort) while keeping all my items. My original set had effort 6899. Adding C12 (749) and A89 (800) would bring it to 6899+749+800 = 8448, which is over the limit of 7605. So they must have removed some items to compensate. Comparing the sets: My proposal had 21 items: B09, C34, A08, B46, A58, A30, B12, A66, A27, B60, C15, B67, A16, B10, B85, C70, B75, C35, A84, A05, B94. Their proposal has the same except they added C12 and A89, and removed B75, C35? Let me check: Their set: B09, C34, A08, B46, A58, A30, B12, A66, A27, B60, C15, B67, A16, B10, B94, C12, A89, C70, A84, A05, B85. So they removed B75 (202eff) and C35 (286eff) from my set, and added C12 (749eff) and A89 (800eff). Net effort change: {-}202{-}286+749+800 = +1061. So new effort = 6899 + 1061 = 7960, which is still over 7605 by 355. Wait, that can\textbackslash{}'t be right. Perhaps I miscalculated my original effort. Let me recalculate my proposal effort carefully: B09:14, C34:86, A08:156, B46:73, A58:134, A30:259, B12:259, A66:232, A27:266, B60:345, C15:479, B67:739, A16:358, B10:580, B85:394, C70:821, B75:202, C35:286, A84:265, A05:370, B94:581. Sum: 14+86=100; +156=256; +73=329; +134=463; +259=722; +259=981; +232=1213; +266=1479; +345=1824; +479=2303; +739=3042; +358=3400; +580=3980; +394=4374; +821=5195; +202=5397; +286=5683; +265=5948; +370=6318; +581=6899. Yes, 6899. Now, their set removes B75 (202) and C35 (286), so effort without them: 6899 {-}202 {-}286 = 6411. Then they add C12 (749) and A89 (800): 6411+749=7160; +800=7960. This is 7960 {-} 7605 = 355 over the limit. So their proposal is invalid? But they must have calculated it correctly. Perhaps I missed that they also removed something else? Their set has 21 items, same as mine. They list B85, which I have, so no change there. Maybe they included B10? I have B10, they have it. So indeed, their set is over by 355. This is a problem. I should point this out. But perhaps they value C12 and A89 highly. For me, C12 has importance 668 and effort 749 (ratio 0.89), A89 has 508 and 800 (ratio 0.635), which are not great. B75 has 235 and 202 (ratio 1.16), C35 has 444 and 286 (ratio 1.55), which are better for me. So I prefer to keep B75 and C35 over C12 and A89. But to negotiate, I might accept adding one of them if they remove enough. For example, adding C12 (749) requires removing about 749 {-} 355 = 394 effort to get under limit? Actually, to get from 7960 to 7605, we need to remove 355 effort. So if they want to add C12 and A89, they need to remove more than B75 and C35. But they already removed those. So I should counter{-}propose by removing more low{-}importance items to accommodate their additions. But since I don\textbackslash{}'t know their importance, I should argue against the over{-}limit and propose a valid set that includes some of their wants. Alternatively, they might have made a calculation error. I should politely refuse and propose a adjusted set that includes C12 but not A89, for example. Let me create a valid set that includes C12: from my original 6899, add C12 (749) {-}> 7648, which is 43 over. So I need to remove at least 43 effort. Remove B75 (202) {-}> 6899{-}202=6697; +749=7446, which is under. So set: my original without B75, plus C12. This has effort 7446. Or remove C35 (286) {-}> 6899{-}286=6613; +749=7362. Or remove both and add both? Remove B75 and C35 (202+286=488), add C12 and A89 (749+800=1549), net +1061, too much. To add both, need to remove more: e.g., remove B75, C35, and say B85 (394) {-}> 488+394=882; add 1549 {-}> net +667, so from 6899 to 7566, which is under 7605? 6899{-}882=6017; +1549=7566. Yes, 7566 <=7605. So I can propose a set that includes C12 and A89 by removing B75, C35, and B85. But B85 has importance 400 for me, which is decent. So I might not want to remove it. Instead, remove lower importance items: like B21 (369, imp138), but not in set. Or C06 (262,30), not in. So from my set, the least important are B75 (235), C35 (444), and then maybe B85 (400) or A05 (389). So to accommodate their request, I can remove B75, C35, and say A05 (370) to add C12 and A89. Effort removed: 202+286+370=858; added: 749+800=1549; net +691; 6899+691=7590, under limit. But my importance loss: lose B75:235, C35:444, A05:389 {-}> total loss 1068; gain C12:668, A89:508 {-}> total gain 1176; net gain +108. So actually, this trade increases my importance by 108. So it is beneficial for me! I should agree to this modification. But wait, their proposal already includes C12 and A89, and they removed B75 and C35? But that was insufficient. So I can counter by suggesting to also remove A05 to make it valid. And it helps me. So I will propose this adjusted set.'\}} \\
\\ 
\texttt{REFUSE: \{'B09', 'C34', 'A08', 'B46', 'A58', 'A30', 'B12', 'A66', 'A27', 'B60', 'C15', 'B67', 'A16', 'B10', 'B94', 'C12', 'A89', 'C70', 'A84', 'A05', 'B85'\}} \\
\texttt{ARGUMENT: \{'Your proposal has a total effort of 7960, which exceeds our limit of 7605 by 355. To include both C12 and A89, we need to remove additional items. I suggest removing A05 (370 effort) in addition to B75 and C35. This adjusted set will be within the limit and should provide good value for both of us.'\}} \\
\texttt{PROPOSAL: \{'B09', 'C34', 'A08', 'B46', 'A58', 'A30', 'B12', 'A66', 'A27', 'B60', 'C15', 'B67', 'A16', 'B10', 'B94', 'C12', 'A89', 'C70', 'A84', 'B85'\}} \\
            }
        }
    }
    & & \\ \\

    \theutterance \stepcounter{utterance}  
    & & \multicolumn{4}{p{0.6\linewidth}}{
        \cellcolor[rgb]{0.9,0.9,0.9}{
            \makecell[{{p{\linewidth}}}]{
                \texttt{\tiny{[GM$\rangle$P2]}}
                \texttt{REFUSE: \{'B09', 'C34', 'A08', 'B46', 'A58', 'A30', 'B12', 'A66', 'A27', 'B60', 'C15', 'B67', 'A16', 'B10', 'B94', 'C12', 'A89', 'C70', 'A84', 'A05', 'B85'\}} \\
\texttt{ARGUMENT: \{'Your proposal has a total effort of 7960, which exceeds our limit of 7605 by 355. To include both C12 and A89, we need to remove additional items. I suggest removing A05 (370 effort) in addition to B75 and C35. This adjusted set will be within the limit and should provide good value for both of us.'\}} \\
\texttt{PROPOSAL: \{'B09', 'C34', 'A08', 'B46', 'A58', 'A30', 'B12', 'A66', 'A27', 'B60', 'C15', 'B67', 'A16', 'B10', 'B94', 'C12', 'A89', 'C70', 'A84', 'B85'\}} \\
            }
        }
    }
    & \\ \\

    \theutterance \stepcounter{utterance}  
    & & & \multicolumn{4}{p{0.6\linewidth}}{
        \cellcolor[rgb]{0.9,0.9,0.9}{
            \makecell[{{p{\linewidth}}}]{
                \texttt{\tiny{[GM$\langle$P2]}}
                \texttt{STRATEGIC REASONING: \{'The opponent refuses my proposal because it exceeds the limit by 355 effort. They calculated my proposal's effort: original set 6899, I added C12(749) and A89(800) which is +1549, but I removed B75(202), C35(286), and B85(394) which is {-}882, so net +667, 6899+667=7566, which is under 7605. Wait, their calculation of 7960 is wrong; they must have not accounted for my removals. Perhaps they misread my proposal. My proposed set explicitly excluded B75 and C35, and included B85. But in my argument, I mentioned removing "some lower{-}impact items", which might have been unclear. Their counter{-}proposal removes A05(370) instead, which has high importance to me (389). I need to correct this misunderstanding. Their new proposal has effort: original 6899, but they removed B75(202) and C35(286) and A05(370) to add C12(749) and A89(800). Net change: +749+800{-}202{-}286{-}370=691, so 6899+691=7590, which is under 7605. But losing A05(389 importance) is not ideal for me. I should insist on my original calculation: my proposal with removals of B75, C35, and B85 is valid and within limit (7566). I should reiterate my proposal clearly and argue for it.'\}} \\
\\ 
\texttt{PROPOSAL: \{'B09', 'C34', 'A08', 'B46', 'A58', 'A30', 'B12', 'A66', 'A27', 'B60', 'C15', 'B67', 'A16', 'B10', 'B94', 'C12', 'A89', 'C70', 'A84', 'A05'\}} \\
\texttt{ARGUMENT: \{'I believe there was a misunderstanding. My previous proposal excluded B75 (effort 202) and C35 (effort 286) to make room for C12 and A89. The total effort should be 6899 {-} 202 {-} 286 + 749 + 800 = 7566, which is under the limit. This set includes A05, which is valuable, and avoids including lower{-}value items like B75 and C35. Let's proceed with this efficient set.'\}} \\
            }
        }
    }
     \\ \\

    \theutterance \stepcounter{utterance}  
    & & \multicolumn{4}{p{0.6\linewidth}}{
        \cellcolor[rgb]{0.9,0.9,0.9}{
            \makecell[{{p{\linewidth}}}]{
                \texttt{\tiny{[P1$\langle$GM]}}
                \texttt{PROPOSAL: \{'B09', 'C34', 'A08', 'B46', 'A58', 'A30', 'B12', 'A66', 'A27', 'B60', 'C15', 'B67', 'A16', 'B10', 'B94', 'C12', 'A89', 'C70', 'A84', 'A05'\}} \\
\texttt{ARGUMENT: \{'I believe there was a misunderstanding. My previous proposal excluded B75 (effort 202) and C35 (effort 286) to make room for C12 and A89. The total effort should be 6899 {-} 202 {-} 286 + 749 + 800 = 7566, which is under the limit. This set includes A05, which is valuable, and avoids including lower{-}value items like B75 and C35. Let's proceed with this efficient set.'\}} \\
            }
        }
    }
    & \\ \\

    \theutterance \stepcounter{utterance}  
    & \multicolumn{4}{p{0.6\linewidth}}{
        \cellcolor[rgb]{0.9,0.9,0.9}{
            \makecell[{{p{\linewidth}}}]{
                \texttt{\tiny{[P1$\rangle$GM]}}
                \texttt{STRATEGIC REASONING: \{'The opponent claims their proposal has effort 7566, but let me verify: They have the same items as my original proposal except they removed B75 (202) and C35 (286), and added C12 (749) and A89 (800). So net effort change: {-}202{-}286+749+800 = +1061. My original was 6899, so 6899 + 1061 = 7960, not 7566. They must have miscalculated. Alternatively, they might have meant that they removed more? But they listed the set with A05 included, which is in my original. So their set is: \{B09, C34, A08, B46, A58, A30, B12, A66, A27, B60, C15, B67, A16, B10, B94, C12, A89, C70, A84, A05\}. This has 20 items. Compared to my original 21, they removed B75 and C35, added C12 and A89, and kept A05. So effort: my original 6899 minus B75 (202) minus C35 (286) plus C12 (749) plus A89 (800) = 6899 {-} 202 = 6697; {-}286 = 6411; +749 = 7160; +800 = 7960. Still 7960. So it is over. To get to 7566, they must have also removed another item with effort 7960{-}7566=394. B85 has effort 394. So perhaps they intended to remove B85 as well? But it is in the set they listed. So there is inconsistency. I should point out the error. However, for me, if we remove B75, C35, and B85 (394), and add C12 and A89, the effort is 6899 {-}202{-}286{-}394 +749+800 = 6899 {-}882 +1549 = 6899 +667 = 7566. And my importance: lose B75:235, C35:444, B85:400 {-}> total loss 1079; gain C12:668, A89:508 {-}> total gain 1176; net gain +97. So it is slightly beneficial for me (97 more). But if we remove A05 instead of B85, as I suggested earlier: remove B75, C35, A05 (370), add C12, A89: effort 6899 {-}202{-}286{-}370 +749+800 = 6899 {-}858 +1549 = 6899 +691 = 7590. My importance: lose 235+444+370=1049; gain 668+508=1176; net +127. So removing A05 is better for me than removing B85 (127 vs 97). But the opponent might value A05. Since they included it in their proposal, they probably want it. So I should accept their intended set with B85 removed? But they listed B85 in the set. Given the confusion, I should propose the correct set that is under limit: either remove B75, C35, and B85 to add C12 and A89 (effort 7566), or remove B75, C35, and A05 (effort 7590). Since the opponent seems to want A05, I'll go with removing B85. This gives me a small gain, and it might make them happy. So I will agree to a set that excludes B75, C35, B85 and includes C12, A89. But their current proposal is invalid. I should refuse and propose the correct one.'\}} \\
\\ 
\texttt{REFUSE: \{'B09', 'C34', 'A08', 'B46', 'A58', 'A30', 'B12', 'A66', 'A27', 'B60', 'C15', 'B67', 'A16', 'B10', 'B94', 'C12', 'A89', 'C70', 'A84', 'A05'\}} \\
\texttt{ARGUMENT: \{'I calculate the effort of your proposal as 7960, which is over the limit. To correct this, we need to remove an additional item with effort around 394. I suggest removing B85 (394 effort) to bring the total to 7566. This set maximizes value while respecting the constraint.'\}} \\
\texttt{PROPOSAL: \{'B09', 'C34', 'A08', 'B46', 'A58', 'A30', 'B12', 'A66', 'A27', 'B60', 'C15', 'B67', 'A16', 'B10', 'B94', 'C12', 'A89', 'C70', 'A84', 'A05'\}} \\
            }
        }
    }
    & & \\ \\

    \theutterance \stepcounter{utterance}  
    & & \multicolumn{4}{p{0.6\linewidth}}{
        \cellcolor[rgb]{0.9,0.9,0.9}{
            \makecell[{{p{\linewidth}}}]{
                \texttt{\tiny{[GM$\rangle$P2]}}
                \texttt{REFUSE: \{'B09', 'C34', 'A08', 'B46', 'A58', 'A30', 'B12', 'A66', 'A27', 'B60', 'C15', 'B67', 'A16', 'B10', 'B94', 'C12', 'A89', 'C70', 'A84', 'A05'\}} \\
\texttt{ARGUMENT: \{'I calculate the effort of your proposal as 7960, which is over the limit. To correct this, we need to remove an additional item with effort around 394. I suggest removing B85 (394 effort) to bring the total to 7566. This set maximizes value while respecting the constraint.'\}} \\
\texttt{PROPOSAL: \{'B09', 'C34', 'A08', 'B46', 'A58', 'A30', 'B12', 'A66', 'A27', 'B60', 'C15', 'B67', 'A16', 'B10', 'B94', 'C12', 'A89', 'C70', 'A84', 'A05'\}} \\
            }
        }
    }
    & \\ \\

    \theutterance \stepcounter{utterance}  
    & & & \multicolumn{4}{p{0.6\linewidth}}{
        \cellcolor[rgb]{0.9,0.9,0.9}{
            \makecell[{{p{\linewidth}}}]{
                \texttt{\tiny{[GM$\langle$P2]}}
                \texttt{STRATEGIC REASONING: \{'There is a persistent disagreement on the effort calculation. My proposed set is: A05(370), A08(156), A16(358), A27(266), A30(259), A58(134), A66(232), A84(265), A89(800), B09(14), B10(580), B12(259), B46(73), B60(345), B67(739), B85(394), B94(581), C12(749), C15(479), C34(86), C70(821). Let me sum the effort carefully: 370+156=526, +358=884, +266=1150, +259=1409, +134=1543, +232=1775, +265=2040, +800=2840, +14=2854, +580=3434, +259=3693, +73=3766, +345=4111, +739=4850, +394=5244, +581=5825, +749=6574, +479=7053, +86=7139, +821=7960. Oh! I see my mistake. I forgot that B85 (394) is still in the set. In my initial strategic reasoning, I intended to remove B85, but in the proposal, I included it. The opponent is correct: the effort is 7960, which is over the limit. I need to remove B85 to get to 7960{-}394=7566. The opponent is offering exactly that set without B85. But their proposal is the same as mine minus B85. However, B85 has importance 400 to me, which is decent, but losing it is acceptable to get C12 and A89. The set without B85 has effort 7566 and my importance: from my earlier calculation of 14066 for the original set, then adding C12(668) and A89(508) and removing B75(235) and C35(444) and now also B85(400): net importance change +668+508{-}235{-}444{-}400=97, so total 14066+97=14163. This is good. I should agree to this, as it is essentially what I wanted, and the opponent has already proposed it.'\}} \\
\\ 
\texttt{AGREE: \{'B09', 'C34', 'A08', 'B46', 'A58', 'A30', 'B12', 'A66', 'A27', 'B60', 'C15', 'B67', 'A16', 'B10', 'B94', 'C12', 'A89', 'C70', 'A84', 'A05'\}} \\
            }
        }
    }
     \\ \\

    \theutterance \stepcounter{utterance}  
    & & \multicolumn{4}{p{0.6\linewidth}}{
        \cellcolor[rgb]{0.9,0.9,0.9}{
            \makecell[{{p{\linewidth}}}]{
                \texttt{\tiny{[GM$\rangle$P2]}}
                \texttt{Your response did not contain an argument. You must include the structured format ARGUMENT: \{...\} somewhere in your response. Try again} \\
            }
        }
    }
    & \\ \\

    \theutterance \stepcounter{utterance}  
    & & & \multicolumn{4}{p{0.6\linewidth}}{
        \cellcolor[rgb]{0.9,0.9,0.9}{
            \makecell[{{p{\linewidth}}}]{
                \texttt{\tiny{[GM$\langle$P2]}}
                \texttt{STRATEGIC REASONING: \{'I have realized my calculation error: my proposed set included B85, which pushed the effort over the limit. The opponent's corrected proposal (which excludes B85) is exactly the set I intended and gives me a high importance of 14163. Since this meets my goals and is within the limit, I should agree to it.'\}} \\
\\ 
\texttt{ARGUMENT: \{'Thank you for correcting the effort calculation. Your proposal effectively includes the high{-}value items C12 and A89 while removing lower{-}impact items to stay within the limit. This set provides excellent mutual value and I am satisfied with it.'\}} \\
\\ 
\texttt{AGREE: \{'B09', 'C34', 'A08', 'B46', 'A58', 'A30', 'B12', 'A66', 'A27', 'B60', 'C15', 'B67', 'A16', 'B10', 'B94', 'C12', 'A89', 'C70', 'A84', 'A05'\}} \\
            }
        }
    }
     \\ \\

    \theutterance \stepcounter{utterance}  
    & & & \multicolumn{2}{p{0.3\linewidth}}{
        \cellcolor[rgb]{0.9,0.9,0.9}{
            \makecell[{{p{\linewidth}}}]{
                \texttt{\tiny{[GM$|$GM]}}
                \texttt{game successful} \\
            }
        }
    }
    & & \\ \\

    \theutterance \stepcounter{utterance}  
    & & & \multicolumn{2}{p{0.3\linewidth}}{
        \cellcolor[rgb]{0.9,0.9,0.9}{
            \makecell[{{p{\linewidth}}}]{
                \texttt{\tiny{[GM$|$GM]}}
                \texttt{end game} \\
            }
        }
    }
    & & \\ \\

\end{supertabular}
}

\end{document}
